\subsection*{Core Ontology Concepts}
\textit{General ontological vocabulary used in the formal framework.}

%-----------------------------------------------------
\subsubsection*{Ontology Types}
%-----------------------------------------------------

\begin{description}[style=nextline,leftmargin=1.5cm]

\item[Ontology]
A formal representation of entities within a domain
and the relationships that hold between them,
expressed in a logic or schema suitable for inference.

\item[Ontology, Applied]
The use of ontological methods to structure and analyze
real-world domains for practical purposes.
Applied ontology is the field;
an applied ontology is an ontology deployed
in a specific domain context.

\item[Ontology, Domain]
An ontology that represents the entities, relations,
and constraints specific to a particular subject area,
such as legal institutions, biomedical processes,
or financial instruments.
Domain ontologies are often built by extending
a foundational or upper ontology with domain-specific vocabulary.
Contrast with \emph{Ontology, Foundational}.

\item[Ontology, Epistemic]
An ontology that represents claims, beliefs,
assertions, reports, or descriptions as objects of representation.
In the setting of this paper,
epistemic modeling allows an ontology to represent
that a causal or normative claim was made,
without asserting the claim's truth as a substrate-layer fact.
Contrast with \emph{Ontology, Metaphysical}.

\item[Ontology, Foundational]
An ontology that provides domain-independent upper-level
categories intended to be extended for specific applications.
Foundational ontologies define general categories
such as continuants, occurrents, roles, processes,
situations, or descriptions,
without committing to a particular domain application.
The Basic Formal Ontology (BFO) is a realist foundational
ontology grounded in scientific practice and distinguishes
continuants from occurrents.
The Descriptive Ontology for Linguistic and Cognitive
Engineering (DOLCE) has a cognitive and linguistic orientation
and supports situation-based modeling of descriptions
and contexts.
Also called \emph{Ontology, Upper}.
Contrast with \emph{Ontology, Domain}.

\item[Ontology, Metaphysical]
An ontology that represents the world directly,
asserting what exists and what relations hold
as ontological facts.
In the terminology of this paper,
metaphysical modeling embeds commitments about
the structure of reality at the substrate layer.
Contrast with \emph{Ontology, Epistemic}.

\item[Ontology, Upper]
See \emph{Ontology, Foundational}.

\end{description}

%-----------------------------------------------------
\subsubsection*{Formal Vocabulary and Assertion}
%-----------------------------------------------------

\begin{description}[style=nextline,leftmargin=1.5cm]

\item[Vocabulary]
The set of symbols available in an ontology for representing
entities, properties, relations, and assertion forms.
In a formal ontology, the vocabulary consists primarily of
predicate symbols together with any constants or functions
used in the representation.
At a given LoA, the vocabulary helps determine
which distinctions can be expressed.
In this paper, the central neutrality constraint concerns
what is asserted as a substrate-layer fact,
not whether a causal or normative symbol can occur
as the content of a reified assertion.

\item[Predicate Symbol]
A symbol in the ontology's formal vocabulary
that denotes a property or relation.
The syntactic occurrence of a predicate symbol
is distinct from the assertion of a proposition using that symbol.
In this paper, neutrality does not require causal or normative
predicate symbols to be impossible or absent from all representation.
It requires that propositions using those predicates
not be asserted as substrate-layer facts when they express
causal or normative commitments.

\item[Ontological Primitive]
A term or predicate that an ontology takes as undefined,
using it to define or constrain other terms
rather than deriving it from prior definitions.
In a neutral substrate, primitives used in substrate-layer
assertions must support framework-invariant referential structure.
Causal and normative predicates may appear as reified content,
but causal and normative commitments may not be asserted
as substrate-layer facts.

\item[Ontological Assertion]
The act of embedding a commitment in an ontology
as a foundational fact.
An ontological assertion is stronger than merely representing
that some agent or framework made a claim:
it treats the asserted proposition as part of the ontology itself.
See \emph{Embedding}.

\end{description}

%-----------------------------------------------------
\subsubsection*{Representation Mechanisms}
%-----------------------------------------------------

\begin{description}[style=nextline,leftmargin=1.5cm]

\item[Representation]
The encoding of information within an ontology.
Representation is broader than embedding:
an ontology may represent a claim via reification
without asserting the claim as a foundational fact.
See also \emph{Embedding} and \emph{Reification}.

\item[Embedding]
The act of asserting a claim as a foundational fact within an ontology,
such that it becomes part of the substrate layer.
Embedded content is treated as true by the ontology itself,
not merely recorded as a claim made by some agent,
document, institution, or framework.
Contrast with \emph{Reification}.

\item[Reification]
The representation of a claim, assertion, description,
or relation as an entity rather than as a direct ontological fact.
Reification allows the ontology to represent discourse
about the world without asserting contested conclusions
about the world itself.
For example, the substrate may represent that a framework
asserted a causal or normative claim,
without embedding that claim as a substrate-layer fact.
See also \emph{Ontology, Epistemic}.
Contrast with \emph{Embedding}.

\end{description}

%-----------------------------------------------------
\subsubsection*{Referential Structure}
%-----------------------------------------------------

\begin{description}[style=nextline,leftmargin=1.5cm]

\item[Individuation]
The component of a referential regime that determines
what counts as one referent.
Individuation specifies the conditions under which
an entity, occurrence, or institutional artifact is treated
as a distinct object of reference.

\item[Co-reference]
The component of a referential regime that determines
when two references denote the same referent.
Co-reference concerns the relation among names, descriptions,
records, identifiers, or representations that point to
a common entity, occurrence, or institutional artifact.

\item[Persistence]
The component of a referential regime that determines
when a referent remains the same across time, records,
jurisdictions, contexts, or transformations.
Persistence concerns the continued identity of the referent
under change.

\item[Referential Regime]
The substrate-level structure by which relevant referents
are individuated, co-referred, and tracked across time or context.
A referential regime comprises three components:
\emph{individuation}, \emph{co-reference}, and \emph{persistence}.
In this paper, referential regimes appear as the scope condition
required for a neutral substrate:
the relevant referential regime must be fixed sufficiently
to remain invariant
across the admissible frameworks under consideration.

\end{description}

%-----------------------------------------------------
\subsubsection*{Architectural Levels}
%-----------------------------------------------------

\begin{description}[style=nextline,leftmargin=1.5cm]

\item[Level of Abstraction (LoA)]
A perspective defined by the observables or predicate symbols
made available for representation at a given level.
In this paper, the foundational LoA constrains what the substrate
may assert as a substrate-layer fact.
Causal and normative predicates may occur as the content
of reified assertions,
but causal and normative commitments may not be asserted
of substrate individuals as substrate-layer facts.

\item[Level of Organization (LoO)]
The architectural layer at which entities, relations,
and assertions are modeled within a system.
Levels of Organization are structural rather than merely scalar:
they differ by direction of dependence
(for example, foundational versus interpretive),
not simply by degree, granularity, or complexity.
A system may have multiple LoOs arranged hierarchically,
with foundational layers supporting higher-order interpretive layers.
In a neutral-substrate architecture,
higher interpretive layers may extend but not revise
substrate-layer commitments.

\item[Substrate]
The foundational layer of an ontology that functions
as a shared representational base across interpretive frameworks.
Formally, the \emph{substrate ontology} is the ontology whose
foundational layer provides entities, occurrences,
institutional artifacts, provenance structures,
and referential regimes.
For brevity, this paper refers to that ontology simply as
the \emph{substrate}.
A neutral substrate does not assert causal or normative commitments
as substrate-layer facts.
It may represent causal or normative claims as reified,
attributed, provenance-bearing content.
Also called \emph{substrate layer} or \emph{foundational layer}.

\end{description}
